
\renewcommand\thesection{\arabic{section}}
\renewcommand{\contentsname}{Cuprins}
\renewcommand{\figurename}{Figura}
\renewcommand{\tablename}{Tabel}

\setcounter{tocdepth}{6}
\setcounter{secnumdepth}{6}


% \usepackage[backend=bibtex,style=numeric,citestyle=numeric,backref=true,sorting=none]{biblatex}




\titleformat{\section}{\normalfont\Large\filleft}{\thesection}{18pt}{}



\definecolor{codegreen}{rgb}{0,0.6,0}
\definecolor{codegray}{rgb}{0.5,0.5,0.5}
\definecolor{codepurple}{rgb}{0.58,0,0.82}
\definecolor{backcolour}{rgb}{0.96,0.96,0.96}

\lstdefinestyle{mystyle}{
	backgroundcolor=\color{backcolour},   
	commentstyle=\color{codegreen},
	keywordstyle=\color{blue},
	numberstyle=\tiny\color{codegray},
	stringstyle=\color{codepurple},
	basicstyle=\fontsize{10}{10}\selectfont\ttfamily,
	breakatwhitespace=false,         
	breaklines=true,                 
	captionpos=b,                    
	keepspaces=true,                 
	numbers=left,                    
	numbersep=2pt,                  
	showspaces=false,                
	showstringspaces=false,
	showtabs=false,                  
	tabsize=2
}

\lstset{style=mystyle}




\setmainfont{UTSans-Regular.otf}


	


\newpage



\newpage
\pagenumbering{arabic}
\renewcommand{\bibname}{Bibliografie}
\renewcommand\lstlistlistingname{\normalsize{CODURI SURSĂ}}
\renewcommand{\lstlistingname}{\normalsize{Secvență de Cod}}
\renewcommand{\listfigurename}{FIGURI}
\renewcommand{\cftloftitlefont}{\normalsize}
\renewcommand{\cftfigfont}{\normalsize}
\renewcommand{\cftfigpagefont}{\normalsize}
\renewcommand{\listtablename}{TABELE}
\renewcommand{\cftlottitlefont}{\normalsize}
\renewcommand{\cfttabfont}{\normalsize}
\renewcommand{\cfttabpagefont}{\normalsize}
 \clearpage

% % trebuie completata manual

\clearpage




 \setcounter{page}{4}
\section{Context și motivație}

\subsection{Context}

În societatea contemporană majoritatea domeniilor cum sunt cele de producție, farmaceutică până la apărare sau alimentar prezența echipamentelor de tip roboți sau brațe robotice a devenit tot mai răspândită, unde se  integrează expertiza umane cu cea a asistenței din partea roboților. Astfel a apărut o nevoie pentru testarea și găsirea mai rapidă a erorilor și defectelor unor astfel de echipamente din partea producătorilor și antrenarea directă pe aceștia poate duce la defectarea lor sau chiar la accidentări, ducând la costuri suplimentare atât din punct de vedere al companiei care a achiziționat robotul cât și din partea producătorului robotului.

O soluție pentru testarea și găsirea eventualelor erori și defecte o reprezintă tehnologia "digital twin". Pentru partea de antrenare a operatorului uman se poate folosi realitatea virtuală pentru a implementa un mediu de antrenament ce implementează un "digital twin", mediul cuprinzând majoritatea scenariilor în care operatorul uman și brațul robotic pot să fie supuse în activitățile la care sunt implicate.
%cuprizând toate funcționalitățile brațului real și nefiind necesară interacțiunea directă cu brațul. 


O interacțiune directă dintre robotul real și "digital twin" poate fi utilă în momentul în care este necesară controlul direct al brațului real sau al robotului în medii periculoase cum sunt cele manevrarea elementelor radioactive sau a fiolelor din laboratoare.

Lucrarea își propune realizarea unui astfel de mediu de antrenament în realitatea virtuală folosind implementarea unui  "digital twin"  pentru un braț robotic cu 4 grade de liberate ce permite ridicarea și mutarea de obiecte folosind o pompă. De asemenea acest sistem urmărește ca procesul de antrenare a utilizatorilor umani, care nu au interacționat niciodată sau rar cu un astfel de sistem să fie cât mai simplă, intuitivă astfel încât timpul petrecut în antrenare să fie cât mai mic și nivelul de pregătire dobândit să fie cât mai mare.


%sunt tot mai prezente dispozitive și echipamente electronice care ne ajută și ne îmbunătățesc stilul de viață în majoritatea privințelor. Astfel au apărut 2 probleme principale:
%aceste echipamente și dispozitive sunt construite din mii și mii de piese și componenete și aceste echipamente și dispozitive a apărut o cerere din ce în ce mai mare. 

%Având în vedere aceste aspecte, cererea și complexitatea, dexteritatea și precizia umamă a ajuns să nu mai fie suficientă. Astfel pentru procesele de fabricație au început să fie folosite aproape în totalitate, roboți și brațe robotice, care fie sunt controlate în mod automat, fie sunt controlate de către un operator uman, apărând nevoia antrenării unui om despre cum să folosească aceste brațe/roboți

%În ultimul timp colaborarea dintre oameni și roboți în mediile industriale a prins tracțiune, în special în procesele de asamblare și dezasamblare de mare precizie. Creșterea complexității mediului de lucru industrial, necesită metodologii de antrenare avansate, prin integrarea expertizei umane cu cea a asistenței din partea roboților. Dar această instruire necesită resurse ridicate și implică riscuri de accidentare.


%Tehnologia VR 

%Această lucrare își propune realizarea unui sistem interactiv de antrenare folosind realitatea virtuală pentru mediul industrial care lucrează cu roboți care să faciliteze antrenarea utilizatorilor umani neexperimentați la un mediul complet de antrenare în folosirea roboților fără implicarea de riscuri.

\subsection{Motivație}
Motivația principală din spatele realizării constă în creșterea în popularitate a tehnologiei de "digital twin" în domenii precum:medical, de producție  domeniul energetic \cite{DigTwin}.
Pe partea de producție aceștia pot ajuta operatorul să înțeleagă condițiile asset-ului și să realizeze planuri predictive pentru mentenanță. Producătorii folosind "digital twins" pentru planificarea zonei de producție, pot să simuleze impactul asupra liniei de producție, ajustând și optimizând pentru cel mai bun flux de lucru. În domeniul auto, producătorii deja folosesc "digital twins" pentru a prototipa și simula modele noi de vehicule. Aceștia pot să ofere posibilitatea pentru vehicule pe bază de software și mașini autonome \cite{DT}.
Pentru domeniul energetic, "digital twins" pot fi folosiți în rețele electrice și tehnologia microgrid, pentru a ajuta operatorii în optimizarea și distribuția de energie. Flexibilitatea acestei tehnologii le oferă operatorilor noi unelte pentru a răspunde într-un mod proactiv la deranjamente, să simuleze situații de urgență și în plănuirea integrării surselor de energie regenerabilă eficientă din punct de vedere al costului.
În medicină, le oferă furnizorilor de servicii medicale posibilitatea îmbunătățirii uneltelor de monitorizarea a pacienților. Prin combinarea datelor de la dispozitive multiple destinate pacienților într-un "digital twin", furnizorii pot să obțină un o vedere mai mare de asamblu asupra sănătății pacientului. De asemenea, pot ajuta administratorii să monitorizeze rata de ocupare și îngrijorările legate de siguranță sau chiar să monitorizeze calitatea aerului pentru a ajuta în prevenirea răspândirii patogenilor \cite{DT}.

Un alt domeniu în care această tehnologie poate fi folosită este cea a telecomunicațiilor. Aceasta le poate oferi furnizorilor din telecom să obține o imagine de ansamblu mai bună asupra  infrastructurii rețelei, a stării asset-ului și să ofere posibilitatea de a se realiza mentenanță predictivă pentru minimizarea momentelor de picare a rețelei. Pe lângă furnizorii îmi mai pot folosi în simularea fluctuării cererilor din rețea pentru teste de stres sau planificarea alternativelor fără necesitate de schimbări costisitoare în lumea reală \cite{DT}.



Un avantaj principal al acestui proiect constă în folosirea tehnologie de realitate virtuală, care reprezintă o soluție modernă care permite simularea și controlul acestor roboți, oferindu-le utilizatorilor posibilitatea de a învăța cum să-i controleze, limitele lor și scenariile de utilizare fără a fi necesar interacțiunea fizică cu cei reali. Un sistem în realitatea virtuală poate reduce costurile, crește eficiența în antrenare și de a îmbunătăți nivelul de pregătire fără a afecta productivitatea reală.

Un alt aspect important al proiectului constă în implementarea a 2 metode de control al brațului din VR. Această metodă de implementare a metodei de control prezintă un avantaj semnificativ deoarece oferă utilizatorului posibilitatea de control al robotului într-un mod simplu și ușor de înțeles, cât și producătorului să testeze rapid implementări și să rezolve eventualele probleme apărute.





\section{Obiectivele lucrării}

Principalele obiective stabilite pentru realizarea lucrării:
\begin{itemize}
	\item Dezvoltatea unui sistem de antrenament care să ofere un minim de pregătire asupra controlului robotului
	\item Realizarea modelului 3D al brațului din realitate
	\item Importarea modelului 3D al brațului în mediul VR din Unity
	\item Aplicarea proprietăților și constrângerilor robotului care sunt necesare pentru mișcarea acestuia
	\item Realizarea primului mod de mișcare al robotului baza pe asociere directă între butoanele de la manete și motoare
	\item Realizarea celui de al 2-lea mod de control al robotului bazat pe cinematică inversă
	\item Aplicarea valorilor citite din Unity asupra robotului din lumea reală pentru a putea fi controlat
	\item Primirea feedback-ului de la robotul real legat de valorile pe care le-a recepționat
\end{itemize}


\clearpage
%Scopul lucrării este de a crea un mediu de simulare interactivă , simplu și sigur pentru antrenarea în folosirea roboților care este destinată utilizatorilor neexperimentați.

%De asemenea acest sistem urmărește ca procesul de antrenare a utilizatorilor umani, care nu au interacționat niciodată sau rar cu un astfel de sistem să fie cât mai simplă, intuitivă astfel încât timpul petrecut în antrenare să fie cât mai mic și nivelul de pregătire dobândit să fie cât mai mare.

%Nu în ultimul rând prin folosirea acestui mediu de antrenare procesul de antrenare nu implică riscurile aferente, dacă antrenarea implica roboții reali.



%Această lucrare prezintă următoarea structură:capitolul 2 prezintă tehnologiile și fundamentele teoretice care stau la baza realizării acestui proiect, capitolul 3 prezintă echipamentele folosite, implementarea, arhitectura aplicației în  realitatea virtuală pentru și modul în care robotul a fost transpus în acest mediu, iar capitolul 4 cuprinde testele realizate și rezultatele în urma testelor, modul în care au fost îndeplinite obiectivele, utilitatea acestui proiect și direcțiile viitoare.
     



% \nocite{*}


